\paragraph{C.4q (Proved: the exact contribution of curves in proper abelian subvarieties).}
Let $A$ be a simple complex abelian variety of dimension
$g\in\{2,3\}$, with
$\operatorname{End}(A)=\mathbb Z$ and
$\operatorname{NS}(A)=\mathbb Z\theta$, where $\theta$ is an
ample primitive class of type $(1,d_2,\ldots,d_g)$.
Put $e=d_g$, and write
\[
 e=u v^2\quad(u\text{ square-free}),\qquad
 \kappa(e)=uv=\prod_{\ell\mid e}\ell^{\lceil v_\ell(e)/2\rceil}.
\]
Let $X=A\times\widehat A$, and let $\mathcal C_{\mathrm{prop}}(X)$
be the subgroup of $H^{4g-2}(X,\mathbb Z)$ generated by the
classes of integral curves contained in translates of proper
abelian subvarieties of $X$. We use the integral cohomological
Fourier transform, followed by the canonical interchange of the
factors of $\widehat X=\widehat A\times A$, to identify
$H^{4g-2}(X,\mathbb Z)$ with $H^2(X,\mathbb Z)$.
If $D$ and $\widehat D$ are the pullbacks of the primitive ample
classes on the two factors and $P=c_1(\mathcal P_A)$, then
\begin{equation}\label{eq:c4-proper-abelian-curve-lattice}
 \boxed{\quad
 \mathcal F_X\bigl(\mathcal C_{\mathrm{prop}}(X)\bigr)
 =\mathbb ZD\oplus\kappa(e)\mathbb ZP
                    \oplus\mathbb Z\widehat D.
 \quad}
\end{equation}
Here the factor interchange is understood in the notation
$\mathcal F_X$. Consequently the smallest positive multiple of
\[
 \mathcal R_A=\frac{P^{2g-1}}{(2g-1)!}
\]
that belongs to $\mathcal C_{\mathrm{prop}}(X)$ is exactly
$\kappa(e)$. In particular, this restricted source of curves
cannot remove any prime dividing $e$ from the direct-summand
problem. This is a statement about a specified subgroup of
curve classes, not a counterexample to algebraicity of
$\mathcal R_A$.

\emph{Proof.}
Write $\lambda=\varphi_\theta:A\to\widehat A$. The primitive
ample generator on $\widehat A$ has homomorphism
$\widehat\lambda=e\lambda^{-1}$; this is integral and primitive
because its elementary divisors are
$(e/d_g,e/d_{g-1},\ldots,e/d_1)$, where $d_1=1$.
The product decomposition of the N\'eron--Severi group and
$\operatorname{Hom}(A,\widehat{\widehat A})=\mathbb Z$
give
\[
 \operatorname{NS}(X)=\mathbb ZD\oplus\mathbb ZP
                         \oplus\mathbb Z\widehat D.
\]
The symmetric homomorphism associated to $aD+bP+c\widehat D$
is
\begin{equation}\label{eq:c4-polarization-block}
 \begin{pmatrix}
  a\lambda&b\,\operatorname{id}_{\widehat A}\\
  b\,\operatorname{id}_A&c\widehat\lambda
 \end{pmatrix}:A\times\widehat A\longrightarrow\widehat A\times A.
\end{equation}
After rationally identifying $\widehat A$ with $A$ by $\lambda$,
its two-by-two scalar matrix has determinant $eac-b^2$.
In particular, a nonzero such homomorphism with image dimension
$g$ satisfies
\begin{equation}\label{eq:c4-rank-one-discriminant}
                             eac=b^2.
\end{equation}

The integral Fourier isomorphism and its preservation of Hodge
type are \citet[Proposition 1, p.~241]{frC4Beauville1983}.
For an inclusion $i:B\hookrightarrow X$ of abelian varieties,
the projection formula and the identity of Poincar\'e bundles
\[
 (i\times\operatorname{id}_{\widehat X})^*\mathcal P_X
   =(\operatorname{id}_B\times\widehat i)^*\mathcal P_B
\]
give directly, on cohomology,
\[
               \mathcal F_X(i_*\alpha)
                   =\widehat i^{,*}\mathcal F_B(\alpha).
\]
This formula is proved by the displayed bundle identity, so we
do not invoke an isogeny-only statement for the inclusion $i$.
Translations act trivially on cohomology.

Every positive-dimensional proper abelian subvariety of $X$
has dimension $g$, since $A$ is simple and $X$ is isogenous to
$A^2$. The Fourier transform of an integral curve class on such
a subvariety $B$ is an integral Hodge class in
$H^2(\widehat B,\mathbb Z)$, hence a divisor class by the
Lefschetz $(1,1)$ theorem. Its pullback to $\widehat X$ has
homomorphism rank at most $g$.
Thus its nonzero block matrix satisfies
\eqref{eq:c4-rank-one-discriminant}. For every prime $\ell\mid e$,
\[
 2v_\ell(b)=v_\ell(e)+v_\ell(a)+v_\ell(c)
                  \ge v_\ell(e)
\]
when $abc\ne0$; if $ac=0$, then $b=0$.
It follows that $\kappa(e)$ divides $b$. This proves one inclusion
in \eqref{eq:c4-proper-abelian-curve-lattice}.

For the converse, we first recall that every integral Hodge
class of curves on any complex abelian variety of dimension
at most three is algebraic. In dimension two this is the
Lefschetz $(1,1)$ theorem. In dimension three it is
\citet[Corollary 3.1.9, pp.~29--30]{frC4Grabowski2004}; its
reduction to principally polarized varieties is Proposition
3.1.8, pp.~28--29. Therefore every integral divisor class on
the dual of such a variety is the Fourier transform of an
integral algebraic curve cycle. In particular, curves in the
two factors of $X$ supply the two diagonal generators $D$ and
$\widehat D$.

Consider the divisor class
\[
                 E=D+\kappa(e)P+u\widehat D.
\]
It is integral, and its associated Hermitian form is positive
semidefinite of rank $g$: both diagonal coefficients are
positive and $e\cdot1\cdot u=\kappa(e)^2$.
Let $q:X\to Y$ be the quotient by the connected kernel of that
form. Then $\dim Y=g$, and $E=q^*\eta$ for an integral ample
class $\eta\in\operatorname{NS}(Y)$. To see the integrality
without dividing by an isogeny degree, write $X=V/\Lambda$.
The alternating form of $E$ is integral on $\Lambda$ and has
rational radical. It descends to an integral positive Riemann
form on the quotient lattice by that radical. The
Appell--Humbert construction gives the asserted class $\eta$.

Choose an integral curve cycle $z$ on $\widehat Y$ whose
Fourier transform is $\eta$. The dual of $q$ is an embedding
$\widehat q:\widehat Y\hookrightarrow\widehat X$.
Fourier compatibility shows that its pushforward, followed by
the factor interchange, has Fourier class $E$ (an irrelevant
overall sign can be absorbed into $z$). Every component is
contained in a proper abelian subvariety. Subtracting the two
diagonal classes produces $\kappa(e)P$. Together with $D$ and
$\widehat D$, this proves the reverse inclusion.
Finally,
$\mathcal F_X(\mathcal R_A)=(-1)^{g+1}P$ under our factor
identification; this is also
\citet[Lemma 3.4, equation (6)]{cyBeckmannFortman2023}.
The exact assertion about multiples of $\mathcal R_A$ follows.
\hfill$\square$

\paragraph{The nearby integral Hodge and Tate theorems.}
The distinction in C.4q matters even though integral curve
classes on the original variety $A$ are understood in these
dimensions. In characteristic zero, Grabowski's theorem applies
to every abelian threefold, with no principal-polarization
hypothesis. Its proof transforms integral divisor classes,
expresses them as differences of ample classes, and principalizes
each ample class; this produces curves on the dual threefold.
It does not assert the algebraicity of the Poincar\'e curve
class on the sixfold $A\times\widehat A$.

Likewise, \citet[Theorem 7.1 and Proposition 7.2, pp.~17--19 of
the opened author version]{frC4Totaro2021} proves the integral
Tate conjecture for abelian threefolds over the separable
closure of a finitely generated field. The formulation uses
$\ell$-adic cohomology for $\ell\ne\operatorname{char}(k)$,
as specified in its Section 1; it does not itself supply the
crystalline statement at the characteristic prime.
Nor does it assert the required integral curve theorem on
$A\times\widehat A$. C.4q identifies an explicit limitation of
trying to obtain that cross class only by pushing curves
forward from its proper abelian subvarieties.
